\section{The minimum of the ambient lattice used in dimension 45}
\label{app:ambient}


Let \(C\) be the extended ternary quadratic-residue \([48,24,15]\) code
\cite{nebe}, and put \(M=\{z\in\mathbb Z^{48}:z\bmod3\in C\}\) and
\(L_0=\{z/\sqrt3:z\in M,\ z^2\equiv0\pmod6\}\). The supplied basis
determines a characteristic numerator \(\chi\in M\) with odd
coordinates, \(\chi^2=264\), and \(\chi\cdot z\equiv z^2\pmod6\). Define

\[
\Lambda=L_0\cup\bigl(\chi/(2\sqrt3)+L_0\bigr).
\]

The characteristic congruence and \(\chi^2/12=22\) give integrality and
evenness; the two index-two constructions give covolume one. The program
reconstructs an integer numerator basis \(B\) and verifies
\(BB^{\mathsf T}=12G\) for the exact Gram matrix used by the 298-vector
certificate.

In \(L_0\), squared norm two or four would require \(z^2\le12\). The
code minimum weight 15 forces \(z\bmod3=0\), whose positive even lattice
norms are at least six. In the other coset, \(y=\chi+2z\) has 48 odd
coordinates, so \(y^2/12\ge4\). Equality requires \(y_i=\pm1\).

Let \(s_i\in\{\pm1\}\) represent \(\chi_i\bmod3\). The finite checks
give \(s\in C\), \(\chi\cdot s=66\), and \(\chi_i s_i\equiv1\pmod6\).
Normalizing coordinates by \(s\) gives a code containing the all-one
word. Its two-symbol enumerator is \(1+94t^{24}+t^{48}\)
\cite[equation (20)]{munemasa-tamura}, so binary-valued codewords have
weights 0, 24 or 48. Writing \(y=s(1-2b)\) gives such a binary word of
weight \(w\), but

\[
\chi\cdot z
=\frac{\chi\cdot s-\chi^2}{2}-\sum_i\chi_i s_i b_i
\equiv-99-w\equiv3\pmod6,
\]

contrary to \(z/\sqrt3\in L_0\). Evenness excludes all other positive
norms below six, while \(\sqrt3(e_1+e_2)\) has squared norm six. The
theta-series and design-moment arguments therefore apply to this
specified lattice.
